\documentclass[11pt,a4paper]{article}
% main (standalone preamble for editing)
% vim: nowrap: spell:
% ══════════════════════════════════════════════════════════════════════════════
% Speed Maths --- shared preamble
% Included by every sheet and answer file via:  % main (standalone preamble for editing)
% vim: nowrap: spell:
% ══════════════════════════════════════════════════════════════════════════════
% Speed Maths --- shared preamble
% Included by every sheet and answer file via:  % main (standalone preamble for editing)
% vim: nowrap: spell:
% ══════════════════════════════════════════════════════════════════════════════
% Speed Maths --- shared preamble
% Included by every sheet and answer file via:  \input{../../shared/preamble}
% Editing THIS file changes the look of every sheet/answer across every pillar.
% ══════════════════════════════════════════════════════════════════════════════

\usepackage[a4paper, top=25mm, bottom=25mm, left=22mm, right=22mm]{geometry}
\usepackage{amsmath, amssymb}
\usepackage{enumitem}
\usepackage{booktabs}
\usepackage{array}
\usepackage{parskip}
\usepackage{microtype}
\usepackage{fancyhdr}
\usepackage{titlesec}
\usepackage{mdframed}
\usepackage{xcolor}
\usepackage[hidelinks]{hyperref}
\usepackage{graphicx}
\usepackage{tikz}
\usepackage{pgfplots}
\pgfplotsset{compat=1.18}

% ── Colours ───────────────────────────────────────────────────────────────────
\definecolor{darkgray}{gray}{0.15}
\definecolor{medgray}{gray}{0.45}
\definecolor{lightgray}{gray}{0.93}
\definecolor{boxgray}{gray}{0.88}

% ── Section formatting ──────────────────────────────────────────────────────────
\titleformat{\section}[block]
  {\normalfont\large\bfseries}{}
  {0pt}{}[\vspace{2pt}\hrule\vspace{6pt}]
\titlespacing*{\section}{0pt}{14pt}{4pt}

% ── List formatting ─────────────────────────────────────────────────────────────
\setlist[enumerate,1]{label=\textbf{\arabic*.}, leftmargin=2em, itemsep=10pt}

% ── Branding constants (edit here, not per-file) ────────────────────────────────
\newcommand{\SpeedExamLine}{\textbf{Speed Maths:} Competition \& Exam Prep}
\newcommand{\SpeedCredit}{Series by \href{https://www.linkedin.com/in/cerealdev/}{David Akinyele-Aje}}

% ── Header / footer ──────────────────────────────────────────────────────────────
% Usage (once, after \input{../../shared/preamble} and before \begin{document}):
%   \SpeedHeader{Algebra}{3}
% First arg  = pillar name shown as "Speed <Pillar> --- Daily Drill"
% Second arg = worksheet number
\pagestyle{fancy}
\newcommand{\SpeedHeader}[2]{%
  \fancyhf{}%
  \renewcommand{\headrulewidth}{0.4pt}%
  \setlength{\headheight}{14pt}%
  \fancyhead[L]{\small\textbf{Speed #1 --- Daily Drill}}%
  \fancyhead[R]{\small Worksheet \##2}%
  \fancyfoot[L]{\small \SpeedExamLine}%
  \fancyfoot[C]{\small\thepage}%
  \fancyfoot[R]{\small \SpeedCredit}%
  \renewcommand{\footrulewidth}{0.4pt}%
}

% ── Title block (used at the top of every sheet AND answer file) ───────────────
% Usage: \SpeedTitleBlock{Daily Algebra Drill \#3}{Jane Doe}
% %% AGENTS: When generating a new sheet, ask the human contributor for the name they want credited in argument 2.
% For answer files, pass the "--- Answers \& Investigations" suffix in the arg.
\newcommand{\SpeedTitleBlock}[2]{%
  \begin{center}
    {\LARGE\bfseries #1}\\[4pt]
    {\normalsize\itshape Sheet by #2}\\[6pt]
    \hrule\vspace{4pt}
    \begin{tabular}{llcc}
      \toprule
      \textbf{Section} & \textbf{Topic} & \textbf{Questions} & \textbf{Target time}\\
      \midrule
      A & Rapid Recognition          & 10 & 2 min 30 sec \\
      B & Manipulation Drills        & 10 & 8 min        \\
      C & Substitution \& Structure  & 8  & 10 min       \\
      D & Challenge                  & 5  & 15 min       \\
      \bottomrule
    \end{tabular}
    \vspace{4pt}
    \hrule
  \end{center}
}

% ── Sheet metadata: topic + the day's new tools ────────────────────────────────
% Usage (immediately after \SpeedTitleBlock, sheets only):
%   \SpeedMeta{Expanding, factorising, and the identity toolkit}
%             {difference of squares, sum/product form, completing the square}
%
% Arg 1 = the sheet's topic in one line. The website lists this instead of
%         "Sheet 03", so write the thing a reader would actually search for.
% Arg 2 = comma-separated, at most five: the tools this sheet introduces that
%         earlier sheets in the pillar did not. These become the website's tags.
%         Leave out anything the reader already met — the tags are meant to say
%         what is new today, not to summarise the sheet.
%
% Both are parsed straight out of this file by tools/build_website.py, so the
% sheet stays the single source of truth and the site cannot drift from it.
%
% This renders NOTHING. It is a declaration for the build script, not a piece of
% the sheet's layout: adding it to a sheet must not change that sheet's PDF, so
% the 25 existing sheets can be annotated without a single one of them being
% reissued. If the toolkit should also appear on the page, that is a separate
% decision about the sheet design — make it deliberately, here, not by leaving
% this macro printing something nobody asked it to.
\newcommand{\SpeedMeta}[2]{}

% ── Answer / Method / Investigate-further boxes (answer files only) ────────────
\newmdenv[
  backgroundcolor=lightgray,
  linecolor=medgray,
  linewidth=0.5pt,
  leftmargin=0pt, rightmargin=0pt,
  innerleftmargin=8pt, innerrightmargin=8pt,
  innertopmargin=5pt, innerbottommargin=5pt,
  skipabove=4pt, skipbelow=4pt
]{ansbox}

\newmdenv[
  backgroundcolor=white,
  linecolor=darkgray,
  linewidth=0.8pt,
  leftline=true, rightline=false, topline=false, bottomline=false,
  leftmargin=0pt, rightmargin=0pt,
  innerleftmargin=8pt, innerrightmargin=6pt,
  innertopmargin=4pt, innerbottommargin=4pt,
  skipabove=4pt, skipbelow=6pt
]{invbox}

\newcommand{\ans}[1]{%
  \begin{ansbox}
    \textbf{Answer:} #1
  \end{ansbox}}

\newcommand{\method}[1]{%
  {\small\textbf{Method:} #1}\par}

\newcommand{\inv}[1]{%
  \begin{invbox}
    \textbf{\small Investigate further:} {\small #1}
  \end{invbox}}

% ── Misc helpers ────────────────────────────────────────────────────────────────
\newcommand{\qn}[1]{\textbf{#1.}\;}

% Mistake-linking: attach a Top 5 Patterns / Common Traps bullet to the question
% label(s) in this sheet where it actually shows up, e.g. \seealso{B6, D2}
\newcommand{\seealso}[1]{\hfill\textit{\footnotesize(see #1)}}

% ── Graph options macro ─────────────────────────────────────────────────────────
% Usage: \GraphOptions{tikz code for A}{tikz code for B}{tikz code for C}{tikz code for D}
\newcommand{\GraphOptions}[4]{%
  \begin{center}
    \begin{tabular}{cc}
      \textbf{A)} & \textbf{B)} \\
      #1 & #2 \\[10pt]
      \textbf{C)} & \textbf{D)} \\
      #3 & #4
    \end{tabular}
  \end{center}
}

% ── Closing motivational rule + cross-promo footer (sheets only) ────────────────
% Usage: \SpeedClosing{"Quote text."}
\newcommand{\SpeedClosing}[1]{%
  \vspace{20pt}
  \begin{center}
  \rule{0.6\textwidth}{0.4pt}\\[4pt]
  {\small\itshape #1}\\[4pt]
  \rule{0.6\textwidth}{0.4pt}
  \end{center}
  \begin{center}
  {\small Found a mistake or want to contribute a sheet?}\\
  {\small Visit the open-source project at \href{https://github.com/The-CerealDev/speed-maths}{github.com/The-CerealDev/speed-maths}}
  \end{center}
}

% Editing THIS file changes the look of every sheet/answer across every pillar.
% ══════════════════════════════════════════════════════════════════════════════

\usepackage[a4paper, top=25mm, bottom=25mm, left=22mm, right=22mm]{geometry}
\usepackage{amsmath, amssymb}
\usepackage{enumitem}
\usepackage{booktabs}
\usepackage{array}
\usepackage{parskip}
\usepackage{microtype}
\usepackage{fancyhdr}
\usepackage{titlesec}
\usepackage{mdframed}
\usepackage{xcolor}
\usepackage[hidelinks]{hyperref}
\usepackage{graphicx}
\usepackage{tikz}
\usepackage{pgfplots}
\pgfplotsset{compat=1.18}

% ── Colours ───────────────────────────────────────────────────────────────────
\definecolor{darkgray}{gray}{0.15}
\definecolor{medgray}{gray}{0.45}
\definecolor{lightgray}{gray}{0.93}
\definecolor{boxgray}{gray}{0.88}

% ── Section formatting ──────────────────────────────────────────────────────────
\titleformat{\section}[block]
  {\normalfont\large\bfseries}{}
  {0pt}{}[\vspace{2pt}\hrule\vspace{6pt}]
\titlespacing*{\section}{0pt}{14pt}{4pt}

% ── List formatting ─────────────────────────────────────────────────────────────
\setlist[enumerate,1]{label=\textbf{\arabic*.}, leftmargin=2em, itemsep=10pt}

% ── Branding constants (edit here, not per-file) ────────────────────────────────
\newcommand{\SpeedExamLine}{\textbf{Speed Maths:} Competition \& Exam Prep}
\newcommand{\SpeedCredit}{Series by \href{https://www.linkedin.com/in/cerealdev/}{David Akinyele-Aje}}

% ── Header / footer ──────────────────────────────────────────────────────────────
% Usage (once, after % main (standalone preamble for editing)
% vim: nowrap: spell:
% ══════════════════════════════════════════════════════════════════════════════
% Speed Maths --- shared preamble
% Included by every sheet and answer file via:  \input{../../shared/preamble}
% Editing THIS file changes the look of every sheet/answer across every pillar.
% ══════════════════════════════════════════════════════════════════════════════

\usepackage[a4paper, top=25mm, bottom=25mm, left=22mm, right=22mm]{geometry}
\usepackage{amsmath, amssymb}
\usepackage{enumitem}
\usepackage{booktabs}
\usepackage{array}
\usepackage{parskip}
\usepackage{microtype}
\usepackage{fancyhdr}
\usepackage{titlesec}
\usepackage{mdframed}
\usepackage{xcolor}
\usepackage[hidelinks]{hyperref}
\usepackage{graphicx}
\usepackage{tikz}
\usepackage{pgfplots}
\pgfplotsset{compat=1.18}

% ── Colours ───────────────────────────────────────────────────────────────────
\definecolor{darkgray}{gray}{0.15}
\definecolor{medgray}{gray}{0.45}
\definecolor{lightgray}{gray}{0.93}
\definecolor{boxgray}{gray}{0.88}

% ── Section formatting ──────────────────────────────────────────────────────────
\titleformat{\section}[block]
  {\normalfont\large\bfseries}{}
  {0pt}{}[\vspace{2pt}\hrule\vspace{6pt}]
\titlespacing*{\section}{0pt}{14pt}{4pt}

% ── List formatting ─────────────────────────────────────────────────────────────
\setlist[enumerate,1]{label=\textbf{\arabic*.}, leftmargin=2em, itemsep=10pt}

% ── Branding constants (edit here, not per-file) ────────────────────────────────
\newcommand{\SpeedExamLine}{\textbf{Speed Maths:} Competition \& Exam Prep}
\newcommand{\SpeedCredit}{Series by \href{https://www.linkedin.com/in/cerealdev/}{David Akinyele-Aje}}

% ── Header / footer ──────────────────────────────────────────────────────────────
% Usage (once, after \input{../../shared/preamble} and before \begin{document}):
%   \SpeedHeader{Algebra}{3}
% First arg  = pillar name shown as "Speed <Pillar> --- Daily Drill"
% Second arg = worksheet number
\pagestyle{fancy}
\newcommand{\SpeedHeader}[2]{%
  \fancyhf{}%
  \renewcommand{\headrulewidth}{0.4pt}%
  \setlength{\headheight}{14pt}%
  \fancyhead[L]{\small\textbf{Speed #1 --- Daily Drill}}%
  \fancyhead[R]{\small Worksheet \##2}%
  \fancyfoot[L]{\small \SpeedExamLine}%
  \fancyfoot[C]{\small\thepage}%
  \fancyfoot[R]{\small \SpeedCredit}%
  \renewcommand{\footrulewidth}{0.4pt}%
}

% ── Title block (used at the top of every sheet AND answer file) ───────────────
% Usage: \SpeedTitleBlock{Daily Algebra Drill \#3}{Jane Doe}
% %% AGENTS: When generating a new sheet, ask the human contributor for the name they want credited in argument 2.
% For answer files, pass the "--- Answers \& Investigations" suffix in the arg.
\newcommand{\SpeedTitleBlock}[2]{%
  \begin{center}
    {\LARGE\bfseries #1}\\[4pt]
    {\normalsize\itshape Sheet by #2}\\[6pt]
    \hrule\vspace{4pt}
    \begin{tabular}{llcc}
      \toprule
      \textbf{Section} & \textbf{Topic} & \textbf{Questions} & \textbf{Target time}\\
      \midrule
      A & Rapid Recognition          & 10 & 2 min 30 sec \\
      B & Manipulation Drills        & 10 & 8 min        \\
      C & Substitution \& Structure  & 8  & 10 min       \\
      D & Challenge                  & 5  & 15 min       \\
      \bottomrule
    \end{tabular}
    \vspace{4pt}
    \hrule
  \end{center}
}

% ── Sheet metadata: topic + the day's new tools ────────────────────────────────
% Usage (immediately after \SpeedTitleBlock, sheets only):
%   \SpeedMeta{Expanding, factorising, and the identity toolkit}
%             {difference of squares, sum/product form, completing the square}
%
% Arg 1 = the sheet's topic in one line. The website lists this instead of
%         "Sheet 03", so write the thing a reader would actually search for.
% Arg 2 = comma-separated, at most five: the tools this sheet introduces that
%         earlier sheets in the pillar did not. These become the website's tags.
%         Leave out anything the reader already met — the tags are meant to say
%         what is new today, not to summarise the sheet.
%
% Both are parsed straight out of this file by tools/build_website.py, so the
% sheet stays the single source of truth and the site cannot drift from it.
%
% This renders NOTHING. It is a declaration for the build script, not a piece of
% the sheet's layout: adding it to a sheet must not change that sheet's PDF, so
% the 25 existing sheets can be annotated without a single one of them being
% reissued. If the toolkit should also appear on the page, that is a separate
% decision about the sheet design — make it deliberately, here, not by leaving
% this macro printing something nobody asked it to.
\newcommand{\SpeedMeta}[2]{}

% ── Answer / Method / Investigate-further boxes (answer files only) ────────────
\newmdenv[
  backgroundcolor=lightgray,
  linecolor=medgray,
  linewidth=0.5pt,
  leftmargin=0pt, rightmargin=0pt,
  innerleftmargin=8pt, innerrightmargin=8pt,
  innertopmargin=5pt, innerbottommargin=5pt,
  skipabove=4pt, skipbelow=4pt
]{ansbox}

\newmdenv[
  backgroundcolor=white,
  linecolor=darkgray,
  linewidth=0.8pt,
  leftline=true, rightline=false, topline=false, bottomline=false,
  leftmargin=0pt, rightmargin=0pt,
  innerleftmargin=8pt, innerrightmargin=6pt,
  innertopmargin=4pt, innerbottommargin=4pt,
  skipabove=4pt, skipbelow=6pt
]{invbox}

\newcommand{\ans}[1]{%
  \begin{ansbox}
    \textbf{Answer:} #1
  \end{ansbox}}

\newcommand{\method}[1]{%
  {\small\textbf{Method:} #1}\par}

\newcommand{\inv}[1]{%
  \begin{invbox}
    \textbf{\small Investigate further:} {\small #1}
  \end{invbox}}

% ── Misc helpers ────────────────────────────────────────────────────────────────
\newcommand{\qn}[1]{\textbf{#1.}\;}

% Mistake-linking: attach a Top 5 Patterns / Common Traps bullet to the question
% label(s) in this sheet where it actually shows up, e.g. \seealso{B6, D2}
\newcommand{\seealso}[1]{\hfill\textit{\footnotesize(see #1)}}

% ── Graph options macro ─────────────────────────────────────────────────────────
% Usage: \GraphOptions{tikz code for A}{tikz code for B}{tikz code for C}{tikz code for D}
\newcommand{\GraphOptions}[4]{%
  \begin{center}
    \begin{tabular}{cc}
      \textbf{A)} & \textbf{B)} \\
      #1 & #2 \\[10pt]
      \textbf{C)} & \textbf{D)} \\
      #3 & #4
    \end{tabular}
  \end{center}
}

% ── Closing motivational rule + cross-promo footer (sheets only) ────────────────
% Usage: \SpeedClosing{"Quote text."}
\newcommand{\SpeedClosing}[1]{%
  \vspace{20pt}
  \begin{center}
  \rule{0.6\textwidth}{0.4pt}\\[4pt]
  {\small\itshape #1}\\[4pt]
  \rule{0.6\textwidth}{0.4pt}
  \end{center}
  \begin{center}
  {\small Found a mistake or want to contribute a sheet?}\\
  {\small Visit the open-source project at \href{https://github.com/The-CerealDev/speed-maths}{github.com/The-CerealDev/speed-maths}}
  \end{center}
}
 and before \begin{document}):
%   \SpeedHeader{Algebra}{3}
% First arg  = pillar name shown as "Speed <Pillar> --- Daily Drill"
% Second arg = worksheet number
\pagestyle{fancy}
\newcommand{\SpeedHeader}[2]{%
  \fancyhf{}%
  \renewcommand{\headrulewidth}{0.4pt}%
  \setlength{\headheight}{14pt}%
  \fancyhead[L]{\small\textbf{Speed #1 --- Daily Drill}}%
  \fancyhead[R]{\small Worksheet \##2}%
  \fancyfoot[L]{\small \SpeedExamLine}%
  \fancyfoot[C]{\small\thepage}%
  \fancyfoot[R]{\small \SpeedCredit}%
  \renewcommand{\footrulewidth}{0.4pt}%
}

% ── Title block (used at the top of every sheet AND answer file) ───────────────
% Usage: \SpeedTitleBlock{Daily Algebra Drill \#3}{Jane Doe}
% %% AGENTS: When generating a new sheet, ask the human contributor for the name they want credited in argument 2.
% For answer files, pass the "--- Answers \& Investigations" suffix in the arg.
\newcommand{\SpeedTitleBlock}[2]{%
  \begin{center}
    {\LARGE\bfseries #1}\\[4pt]
    {\normalsize\itshape Sheet by #2}\\[6pt]
    \hrule\vspace{4pt}
    \begin{tabular}{llcc}
      \toprule
      \textbf{Section} & \textbf{Topic} & \textbf{Questions} & \textbf{Target time}\\
      \midrule
      A & Rapid Recognition          & 10 & 2 min 30 sec \\
      B & Manipulation Drills        & 10 & 8 min        \\
      C & Substitution \& Structure  & 8  & 10 min       \\
      D & Challenge                  & 5  & 15 min       \\
      \bottomrule
    \end{tabular}
    \vspace{4pt}
    \hrule
  \end{center}
}

% ── Sheet metadata: topic + the day's new tools ────────────────────────────────
% Usage (immediately after \SpeedTitleBlock, sheets only):
%   \SpeedMeta{Expanding, factorising, and the identity toolkit}
%             {difference of squares, sum/product form, completing the square}
%
% Arg 1 = the sheet's topic in one line. The website lists this instead of
%         "Sheet 03", so write the thing a reader would actually search for.
% Arg 2 = comma-separated, at most five: the tools this sheet introduces that
%         earlier sheets in the pillar did not. These become the website's tags.
%         Leave out anything the reader already met — the tags are meant to say
%         what is new today, not to summarise the sheet.
%
% Both are parsed straight out of this file by tools/build_website.py, so the
% sheet stays the single source of truth and the site cannot drift from it.
%
% This renders NOTHING. It is a declaration for the build script, not a piece of
% the sheet's layout: adding it to a sheet must not change that sheet's PDF, so
% the 25 existing sheets can be annotated without a single one of them being
% reissued. If the toolkit should also appear on the page, that is a separate
% decision about the sheet design — make it deliberately, here, not by leaving
% this macro printing something nobody asked it to.
\newcommand{\SpeedMeta}[2]{}

% ── Answer / Method / Investigate-further boxes (answer files only) ────────────
\newmdenv[
  backgroundcolor=lightgray,
  linecolor=medgray,
  linewidth=0.5pt,
  leftmargin=0pt, rightmargin=0pt,
  innerleftmargin=8pt, innerrightmargin=8pt,
  innertopmargin=5pt, innerbottommargin=5pt,
  skipabove=4pt, skipbelow=4pt
]{ansbox}

\newmdenv[
  backgroundcolor=white,
  linecolor=darkgray,
  linewidth=0.8pt,
  leftline=true, rightline=false, topline=false, bottomline=false,
  leftmargin=0pt, rightmargin=0pt,
  innerleftmargin=8pt, innerrightmargin=6pt,
  innertopmargin=4pt, innerbottommargin=4pt,
  skipabove=4pt, skipbelow=6pt
]{invbox}

\newcommand{\ans}[1]{%
  \begin{ansbox}
    \textbf{Answer:} #1
  \end{ansbox}}

\newcommand{\method}[1]{%
  {\small\textbf{Method:} #1}\par}

\newcommand{\inv}[1]{%
  \begin{invbox}
    \textbf{\small Investigate further:} {\small #1}
  \end{invbox}}

% ── Misc helpers ────────────────────────────────────────────────────────────────
\newcommand{\qn}[1]{\textbf{#1.}\;}

% Mistake-linking: attach a Top 5 Patterns / Common Traps bullet to the question
% label(s) in this sheet where it actually shows up, e.g. \seealso{B6, D2}
\newcommand{\seealso}[1]{\hfill\textit{\footnotesize(see #1)}}

% ── Graph options macro ─────────────────────────────────────────────────────────
% Usage: \GraphOptions{tikz code for A}{tikz code for B}{tikz code for C}{tikz code for D}
\newcommand{\GraphOptions}[4]{%
  \begin{center}
    \begin{tabular}{cc}
      \textbf{A)} & \textbf{B)} \\
      #1 & #2 \\[10pt]
      \textbf{C)} & \textbf{D)} \\
      #3 & #4
    \end{tabular}
  \end{center}
}

% ── Closing motivational rule + cross-promo footer (sheets only) ────────────────
% Usage: \SpeedClosing{"Quote text."}
\newcommand{\SpeedClosing}[1]{%
  \vspace{20pt}
  \begin{center}
  \rule{0.6\textwidth}{0.4pt}\\[4pt]
  {\small\itshape #1}\\[4pt]
  \rule{0.6\textwidth}{0.4pt}
  \end{center}
  \begin{center}
  {\small Found a mistake or want to contribute a sheet?}\\
  {\small Visit the open-source project at \href{https://github.com/The-CerealDev/speed-maths}{github.com/The-CerealDev/speed-maths}}
  \end{center}
}

% Editing THIS file changes the look of every sheet/answer across every pillar.
% ══════════════════════════════════════════════════════════════════════════════

\usepackage[a4paper, top=25mm, bottom=25mm, left=22mm, right=22mm]{geometry}
\usepackage{amsmath, amssymb}
\usepackage{enumitem}
\usepackage{booktabs}
\usepackage{array}
\usepackage{parskip}
\usepackage{microtype}
\usepackage{fancyhdr}
\usepackage{titlesec}
\usepackage{mdframed}
\usepackage{xcolor}
\usepackage[hidelinks]{hyperref}
\usepackage{graphicx}
\usepackage{tikz}
\usepackage{pgfplots}
\pgfplotsset{compat=1.18}

% ── Colours ───────────────────────────────────────────────────────────────────
\definecolor{darkgray}{gray}{0.15}
\definecolor{medgray}{gray}{0.45}
\definecolor{lightgray}{gray}{0.93}
\definecolor{boxgray}{gray}{0.88}

% ── Section formatting ──────────────────────────────────────────────────────────
\titleformat{\section}[block]
  {\normalfont\large\bfseries}{}
  {0pt}{}[\vspace{2pt}\hrule\vspace{6pt}]
\titlespacing*{\section}{0pt}{14pt}{4pt}

% ── List formatting ─────────────────────────────────────────────────────────────
\setlist[enumerate,1]{label=\textbf{\arabic*.}, leftmargin=2em, itemsep=10pt}

% ── Branding constants (edit here, not per-file) ────────────────────────────────
\newcommand{\SpeedExamLine}{\textbf{Speed Maths:} Competition \& Exam Prep}
\newcommand{\SpeedCredit}{Series by \href{https://www.linkedin.com/in/cerealdev/}{David Akinyele-Aje}}

% ── Header / footer ──────────────────────────────────────────────────────────────
% Usage (once, after % main (standalone preamble for editing)
% vim: nowrap: spell:
% ══════════════════════════════════════════════════════════════════════════════
% Speed Maths --- shared preamble
% Included by every sheet and answer file via:  % main (standalone preamble for editing)
% vim: nowrap: spell:
% ══════════════════════════════════════════════════════════════════════════════
% Speed Maths --- shared preamble
% Included by every sheet and answer file via:  \input{../../shared/preamble}
% Editing THIS file changes the look of every sheet/answer across every pillar.
% ══════════════════════════════════════════════════════════════════════════════

\usepackage[a4paper, top=25mm, bottom=25mm, left=22mm, right=22mm]{geometry}
\usepackage{amsmath, amssymb}
\usepackage{enumitem}
\usepackage{booktabs}
\usepackage{array}
\usepackage{parskip}
\usepackage{microtype}
\usepackage{fancyhdr}
\usepackage{titlesec}
\usepackage{mdframed}
\usepackage{xcolor}
\usepackage[hidelinks]{hyperref}
\usepackage{graphicx}
\usepackage{tikz}
\usepackage{pgfplots}
\pgfplotsset{compat=1.18}

% ── Colours ───────────────────────────────────────────────────────────────────
\definecolor{darkgray}{gray}{0.15}
\definecolor{medgray}{gray}{0.45}
\definecolor{lightgray}{gray}{0.93}
\definecolor{boxgray}{gray}{0.88}

% ── Section formatting ──────────────────────────────────────────────────────────
\titleformat{\section}[block]
  {\normalfont\large\bfseries}{}
  {0pt}{}[\vspace{2pt}\hrule\vspace{6pt}]
\titlespacing*{\section}{0pt}{14pt}{4pt}

% ── List formatting ─────────────────────────────────────────────────────────────
\setlist[enumerate,1]{label=\textbf{\arabic*.}, leftmargin=2em, itemsep=10pt}

% ── Branding constants (edit here, not per-file) ────────────────────────────────
\newcommand{\SpeedExamLine}{\textbf{Speed Maths:} Competition \& Exam Prep}
\newcommand{\SpeedCredit}{Series by \href{https://www.linkedin.com/in/cerealdev/}{David Akinyele-Aje}}

% ── Header / footer ──────────────────────────────────────────────────────────────
% Usage (once, after \input{../../shared/preamble} and before \begin{document}):
%   \SpeedHeader{Algebra}{3}
% First arg  = pillar name shown as "Speed <Pillar> --- Daily Drill"
% Second arg = worksheet number
\pagestyle{fancy}
\newcommand{\SpeedHeader}[2]{%
  \fancyhf{}%
  \renewcommand{\headrulewidth}{0.4pt}%
  \setlength{\headheight}{14pt}%
  \fancyhead[L]{\small\textbf{Speed #1 --- Daily Drill}}%
  \fancyhead[R]{\small Worksheet \##2}%
  \fancyfoot[L]{\small \SpeedExamLine}%
  \fancyfoot[C]{\small\thepage}%
  \fancyfoot[R]{\small \SpeedCredit}%
  \renewcommand{\footrulewidth}{0.4pt}%
}

% ── Title block (used at the top of every sheet AND answer file) ───────────────
% Usage: \SpeedTitleBlock{Daily Algebra Drill \#3}{Jane Doe}
% %% AGENTS: When generating a new sheet, ask the human contributor for the name they want credited in argument 2.
% For answer files, pass the "--- Answers \& Investigations" suffix in the arg.
\newcommand{\SpeedTitleBlock}[2]{%
  \begin{center}
    {\LARGE\bfseries #1}\\[4pt]
    {\normalsize\itshape Sheet by #2}\\[6pt]
    \hrule\vspace{4pt}
    \begin{tabular}{llcc}
      \toprule
      \textbf{Section} & \textbf{Topic} & \textbf{Questions} & \textbf{Target time}\\
      \midrule
      A & Rapid Recognition          & 10 & 2 min 30 sec \\
      B & Manipulation Drills        & 10 & 8 min        \\
      C & Substitution \& Structure  & 8  & 10 min       \\
      D & Challenge                  & 5  & 15 min       \\
      \bottomrule
    \end{tabular}
    \vspace{4pt}
    \hrule
  \end{center}
}

% ── Sheet metadata: topic + the day's new tools ────────────────────────────────
% Usage (immediately after \SpeedTitleBlock, sheets only):
%   \SpeedMeta{Expanding, factorising, and the identity toolkit}
%             {difference of squares, sum/product form, completing the square}
%
% Arg 1 = the sheet's topic in one line. The website lists this instead of
%         "Sheet 03", so write the thing a reader would actually search for.
% Arg 2 = comma-separated, at most five: the tools this sheet introduces that
%         earlier sheets in the pillar did not. These become the website's tags.
%         Leave out anything the reader already met — the tags are meant to say
%         what is new today, not to summarise the sheet.
%
% Both are parsed straight out of this file by tools/build_website.py, so the
% sheet stays the single source of truth and the site cannot drift from it.
%
% This renders NOTHING. It is a declaration for the build script, not a piece of
% the sheet's layout: adding it to a sheet must not change that sheet's PDF, so
% the 25 existing sheets can be annotated without a single one of them being
% reissued. If the toolkit should also appear on the page, that is a separate
% decision about the sheet design — make it deliberately, here, not by leaving
% this macro printing something nobody asked it to.
\newcommand{\SpeedMeta}[2]{}

% ── Answer / Method / Investigate-further boxes (answer files only) ────────────
\newmdenv[
  backgroundcolor=lightgray,
  linecolor=medgray,
  linewidth=0.5pt,
  leftmargin=0pt, rightmargin=0pt,
  innerleftmargin=8pt, innerrightmargin=8pt,
  innertopmargin=5pt, innerbottommargin=5pt,
  skipabove=4pt, skipbelow=4pt
]{ansbox}

\newmdenv[
  backgroundcolor=white,
  linecolor=darkgray,
  linewidth=0.8pt,
  leftline=true, rightline=false, topline=false, bottomline=false,
  leftmargin=0pt, rightmargin=0pt,
  innerleftmargin=8pt, innerrightmargin=6pt,
  innertopmargin=4pt, innerbottommargin=4pt,
  skipabove=4pt, skipbelow=6pt
]{invbox}

\newcommand{\ans}[1]{%
  \begin{ansbox}
    \textbf{Answer:} #1
  \end{ansbox}}

\newcommand{\method}[1]{%
  {\small\textbf{Method:} #1}\par}

\newcommand{\inv}[1]{%
  \begin{invbox}
    \textbf{\small Investigate further:} {\small #1}
  \end{invbox}}

% ── Misc helpers ────────────────────────────────────────────────────────────────
\newcommand{\qn}[1]{\textbf{#1.}\;}

% Mistake-linking: attach a Top 5 Patterns / Common Traps bullet to the question
% label(s) in this sheet where it actually shows up, e.g. \seealso{B6, D2}
\newcommand{\seealso}[1]{\hfill\textit{\footnotesize(see #1)}}

% ── Graph options macro ─────────────────────────────────────────────────────────
% Usage: \GraphOptions{tikz code for A}{tikz code for B}{tikz code for C}{tikz code for D}
\newcommand{\GraphOptions}[4]{%
  \begin{center}
    \begin{tabular}{cc}
      \textbf{A)} & \textbf{B)} \\
      #1 & #2 \\[10pt]
      \textbf{C)} & \textbf{D)} \\
      #3 & #4
    \end{tabular}
  \end{center}
}

% ── Closing motivational rule + cross-promo footer (sheets only) ────────────────
% Usage: \SpeedClosing{"Quote text."}
\newcommand{\SpeedClosing}[1]{%
  \vspace{20pt}
  \begin{center}
  \rule{0.6\textwidth}{0.4pt}\\[4pt]
  {\small\itshape #1}\\[4pt]
  \rule{0.6\textwidth}{0.4pt}
  \end{center}
  \begin{center}
  {\small Found a mistake or want to contribute a sheet?}\\
  {\small Visit the open-source project at \href{https://github.com/The-CerealDev/speed-maths}{github.com/The-CerealDev/speed-maths}}
  \end{center}
}

% Editing THIS file changes the look of every sheet/answer across every pillar.
% ══════════════════════════════════════════════════════════════════════════════

\usepackage[a4paper, top=25mm, bottom=25mm, left=22mm, right=22mm]{geometry}
\usepackage{amsmath, amssymb}
\usepackage{enumitem}
\usepackage{booktabs}
\usepackage{array}
\usepackage{parskip}
\usepackage{microtype}
\usepackage{fancyhdr}
\usepackage{titlesec}
\usepackage{mdframed}
\usepackage{xcolor}
\usepackage[hidelinks]{hyperref}
\usepackage{graphicx}
\usepackage{tikz}
\usepackage{pgfplots}
\pgfplotsset{compat=1.18}

% ── Colours ───────────────────────────────────────────────────────────────────
\definecolor{darkgray}{gray}{0.15}
\definecolor{medgray}{gray}{0.45}
\definecolor{lightgray}{gray}{0.93}
\definecolor{boxgray}{gray}{0.88}

% ── Section formatting ──────────────────────────────────────────────────────────
\titleformat{\section}[block]
  {\normalfont\large\bfseries}{}
  {0pt}{}[\vspace{2pt}\hrule\vspace{6pt}]
\titlespacing*{\section}{0pt}{14pt}{4pt}

% ── List formatting ─────────────────────────────────────────────────────────────
\setlist[enumerate,1]{label=\textbf{\arabic*.}, leftmargin=2em, itemsep=10pt}

% ── Branding constants (edit here, not per-file) ────────────────────────────────
\newcommand{\SpeedExamLine}{\textbf{Speed Maths:} Competition \& Exam Prep}
\newcommand{\SpeedCredit}{Series by \href{https://www.linkedin.com/in/cerealdev/}{David Akinyele-Aje}}

% ── Header / footer ──────────────────────────────────────────────────────────────
% Usage (once, after % main (standalone preamble for editing)
% vim: nowrap: spell:
% ══════════════════════════════════════════════════════════════════════════════
% Speed Maths --- shared preamble
% Included by every sheet and answer file via:  \input{../../shared/preamble}
% Editing THIS file changes the look of every sheet/answer across every pillar.
% ══════════════════════════════════════════════════════════════════════════════

\usepackage[a4paper, top=25mm, bottom=25mm, left=22mm, right=22mm]{geometry}
\usepackage{amsmath, amssymb}
\usepackage{enumitem}
\usepackage{booktabs}
\usepackage{array}
\usepackage{parskip}
\usepackage{microtype}
\usepackage{fancyhdr}
\usepackage{titlesec}
\usepackage{mdframed}
\usepackage{xcolor}
\usepackage[hidelinks]{hyperref}
\usepackage{graphicx}
\usepackage{tikz}
\usepackage{pgfplots}
\pgfplotsset{compat=1.18}

% ── Colours ───────────────────────────────────────────────────────────────────
\definecolor{darkgray}{gray}{0.15}
\definecolor{medgray}{gray}{0.45}
\definecolor{lightgray}{gray}{0.93}
\definecolor{boxgray}{gray}{0.88}

% ── Section formatting ──────────────────────────────────────────────────────────
\titleformat{\section}[block]
  {\normalfont\large\bfseries}{}
  {0pt}{}[\vspace{2pt}\hrule\vspace{6pt}]
\titlespacing*{\section}{0pt}{14pt}{4pt}

% ── List formatting ─────────────────────────────────────────────────────────────
\setlist[enumerate,1]{label=\textbf{\arabic*.}, leftmargin=2em, itemsep=10pt}

% ── Branding constants (edit here, not per-file) ────────────────────────────────
\newcommand{\SpeedExamLine}{\textbf{Speed Maths:} Competition \& Exam Prep}
\newcommand{\SpeedCredit}{Series by \href{https://www.linkedin.com/in/cerealdev/}{David Akinyele-Aje}}

% ── Header / footer ──────────────────────────────────────────────────────────────
% Usage (once, after \input{../../shared/preamble} and before \begin{document}):
%   \SpeedHeader{Algebra}{3}
% First arg  = pillar name shown as "Speed <Pillar> --- Daily Drill"
% Second arg = worksheet number
\pagestyle{fancy}
\newcommand{\SpeedHeader}[2]{%
  \fancyhf{}%
  \renewcommand{\headrulewidth}{0.4pt}%
  \setlength{\headheight}{14pt}%
  \fancyhead[L]{\small\textbf{Speed #1 --- Daily Drill}}%
  \fancyhead[R]{\small Worksheet \##2}%
  \fancyfoot[L]{\small \SpeedExamLine}%
  \fancyfoot[C]{\small\thepage}%
  \fancyfoot[R]{\small \SpeedCredit}%
  \renewcommand{\footrulewidth}{0.4pt}%
}

% ── Title block (used at the top of every sheet AND answer file) ───────────────
% Usage: \SpeedTitleBlock{Daily Algebra Drill \#3}{Jane Doe}
% %% AGENTS: When generating a new sheet, ask the human contributor for the name they want credited in argument 2.
% For answer files, pass the "--- Answers \& Investigations" suffix in the arg.
\newcommand{\SpeedTitleBlock}[2]{%
  \begin{center}
    {\LARGE\bfseries #1}\\[4pt]
    {\normalsize\itshape Sheet by #2}\\[6pt]
    \hrule\vspace{4pt}
    \begin{tabular}{llcc}
      \toprule
      \textbf{Section} & \textbf{Topic} & \textbf{Questions} & \textbf{Target time}\\
      \midrule
      A & Rapid Recognition          & 10 & 2 min 30 sec \\
      B & Manipulation Drills        & 10 & 8 min        \\
      C & Substitution \& Structure  & 8  & 10 min       \\
      D & Challenge                  & 5  & 15 min       \\
      \bottomrule
    \end{tabular}
    \vspace{4pt}
    \hrule
  \end{center}
}

% ── Sheet metadata: topic + the day's new tools ────────────────────────────────
% Usage (immediately after \SpeedTitleBlock, sheets only):
%   \SpeedMeta{Expanding, factorising, and the identity toolkit}
%             {difference of squares, sum/product form, completing the square}
%
% Arg 1 = the sheet's topic in one line. The website lists this instead of
%         "Sheet 03", so write the thing a reader would actually search for.
% Arg 2 = comma-separated, at most five: the tools this sheet introduces that
%         earlier sheets in the pillar did not. These become the website's tags.
%         Leave out anything the reader already met — the tags are meant to say
%         what is new today, not to summarise the sheet.
%
% Both are parsed straight out of this file by tools/build_website.py, so the
% sheet stays the single source of truth and the site cannot drift from it.
%
% This renders NOTHING. It is a declaration for the build script, not a piece of
% the sheet's layout: adding it to a sheet must not change that sheet's PDF, so
% the 25 existing sheets can be annotated without a single one of them being
% reissued. If the toolkit should also appear on the page, that is a separate
% decision about the sheet design — make it deliberately, here, not by leaving
% this macro printing something nobody asked it to.
\newcommand{\SpeedMeta}[2]{}

% ── Answer / Method / Investigate-further boxes (answer files only) ────────────
\newmdenv[
  backgroundcolor=lightgray,
  linecolor=medgray,
  linewidth=0.5pt,
  leftmargin=0pt, rightmargin=0pt,
  innerleftmargin=8pt, innerrightmargin=8pt,
  innertopmargin=5pt, innerbottommargin=5pt,
  skipabove=4pt, skipbelow=4pt
]{ansbox}

\newmdenv[
  backgroundcolor=white,
  linecolor=darkgray,
  linewidth=0.8pt,
  leftline=true, rightline=false, topline=false, bottomline=false,
  leftmargin=0pt, rightmargin=0pt,
  innerleftmargin=8pt, innerrightmargin=6pt,
  innertopmargin=4pt, innerbottommargin=4pt,
  skipabove=4pt, skipbelow=6pt
]{invbox}

\newcommand{\ans}[1]{%
  \begin{ansbox}
    \textbf{Answer:} #1
  \end{ansbox}}

\newcommand{\method}[1]{%
  {\small\textbf{Method:} #1}\par}

\newcommand{\inv}[1]{%
  \begin{invbox}
    \textbf{\small Investigate further:} {\small #1}
  \end{invbox}}

% ── Misc helpers ────────────────────────────────────────────────────────────────
\newcommand{\qn}[1]{\textbf{#1.}\;}

% Mistake-linking: attach a Top 5 Patterns / Common Traps bullet to the question
% label(s) in this sheet where it actually shows up, e.g. \seealso{B6, D2}
\newcommand{\seealso}[1]{\hfill\textit{\footnotesize(see #1)}}

% ── Graph options macro ─────────────────────────────────────────────────────────
% Usage: \GraphOptions{tikz code for A}{tikz code for B}{tikz code for C}{tikz code for D}
\newcommand{\GraphOptions}[4]{%
  \begin{center}
    \begin{tabular}{cc}
      \textbf{A)} & \textbf{B)} \\
      #1 & #2 \\[10pt]
      \textbf{C)} & \textbf{D)} \\
      #3 & #4
    \end{tabular}
  \end{center}
}

% ── Closing motivational rule + cross-promo footer (sheets only) ────────────────
% Usage: \SpeedClosing{"Quote text."}
\newcommand{\SpeedClosing}[1]{%
  \vspace{20pt}
  \begin{center}
  \rule{0.6\textwidth}{0.4pt}\\[4pt]
  {\small\itshape #1}\\[4pt]
  \rule{0.6\textwidth}{0.4pt}
  \end{center}
  \begin{center}
  {\small Found a mistake or want to contribute a sheet?}\\
  {\small Visit the open-source project at \href{https://github.com/The-CerealDev/speed-maths}{github.com/The-CerealDev/speed-maths}}
  \end{center}
}
 and before \begin{document}):
%   \SpeedHeader{Algebra}{3}
% First arg  = pillar name shown as "Speed <Pillar> --- Daily Drill"
% Second arg = worksheet number
\pagestyle{fancy}
\newcommand{\SpeedHeader}[2]{%
  \fancyhf{}%
  \renewcommand{\headrulewidth}{0.4pt}%
  \setlength{\headheight}{14pt}%
  \fancyhead[L]{\small\textbf{Speed #1 --- Daily Drill}}%
  \fancyhead[R]{\small Worksheet \##2}%
  \fancyfoot[L]{\small \SpeedExamLine}%
  \fancyfoot[C]{\small\thepage}%
  \fancyfoot[R]{\small \SpeedCredit}%
  \renewcommand{\footrulewidth}{0.4pt}%
}

% ── Title block (used at the top of every sheet AND answer file) ───────────────
% Usage: \SpeedTitleBlock{Daily Algebra Drill \#3}{Jane Doe}
% %% AGENTS: When generating a new sheet, ask the human contributor for the name they want credited in argument 2.
% For answer files, pass the "--- Answers \& Investigations" suffix in the arg.
\newcommand{\SpeedTitleBlock}[2]{%
  \begin{center}
    {\LARGE\bfseries #1}\\[4pt]
    {\normalsize\itshape Sheet by #2}\\[6pt]
    \hrule\vspace{4pt}
    \begin{tabular}{llcc}
      \toprule
      \textbf{Section} & \textbf{Topic} & \textbf{Questions} & \textbf{Target time}\\
      \midrule
      A & Rapid Recognition          & 10 & 2 min 30 sec \\
      B & Manipulation Drills        & 10 & 8 min        \\
      C & Substitution \& Structure  & 8  & 10 min       \\
      D & Challenge                  & 5  & 15 min       \\
      \bottomrule
    \end{tabular}
    \vspace{4pt}
    \hrule
  \end{center}
}

% ── Sheet metadata: topic + the day's new tools ────────────────────────────────
% Usage (immediately after \SpeedTitleBlock, sheets only):
%   \SpeedMeta{Expanding, factorising, and the identity toolkit}
%             {difference of squares, sum/product form, completing the square}
%
% Arg 1 = the sheet's topic in one line. The website lists this instead of
%         "Sheet 03", so write the thing a reader would actually search for.
% Arg 2 = comma-separated, at most five: the tools this sheet introduces that
%         earlier sheets in the pillar did not. These become the website's tags.
%         Leave out anything the reader already met — the tags are meant to say
%         what is new today, not to summarise the sheet.
%
% Both are parsed straight out of this file by tools/build_website.py, so the
% sheet stays the single source of truth and the site cannot drift from it.
%
% This renders NOTHING. It is a declaration for the build script, not a piece of
% the sheet's layout: adding it to a sheet must not change that sheet's PDF, so
% the 25 existing sheets can be annotated without a single one of them being
% reissued. If the toolkit should also appear on the page, that is a separate
% decision about the sheet design — make it deliberately, here, not by leaving
% this macro printing something nobody asked it to.
\newcommand{\SpeedMeta}[2]{}

% ── Answer / Method / Investigate-further boxes (answer files only) ────────────
\newmdenv[
  backgroundcolor=lightgray,
  linecolor=medgray,
  linewidth=0.5pt,
  leftmargin=0pt, rightmargin=0pt,
  innerleftmargin=8pt, innerrightmargin=8pt,
  innertopmargin=5pt, innerbottommargin=5pt,
  skipabove=4pt, skipbelow=4pt
]{ansbox}

\newmdenv[
  backgroundcolor=white,
  linecolor=darkgray,
  linewidth=0.8pt,
  leftline=true, rightline=false, topline=false, bottomline=false,
  leftmargin=0pt, rightmargin=0pt,
  innerleftmargin=8pt, innerrightmargin=6pt,
  innertopmargin=4pt, innerbottommargin=4pt,
  skipabove=4pt, skipbelow=6pt
]{invbox}

\newcommand{\ans}[1]{%
  \begin{ansbox}
    \textbf{Answer:} #1
  \end{ansbox}}

\newcommand{\method}[1]{%
  {\small\textbf{Method:} #1}\par}

\newcommand{\inv}[1]{%
  \begin{invbox}
    \textbf{\small Investigate further:} {\small #1}
  \end{invbox}}

% ── Misc helpers ────────────────────────────────────────────────────────────────
\newcommand{\qn}[1]{\textbf{#1.}\;}

% Mistake-linking: attach a Top 5 Patterns / Common Traps bullet to the question
% label(s) in this sheet where it actually shows up, e.g. \seealso{B6, D2}
\newcommand{\seealso}[1]{\hfill\textit{\footnotesize(see #1)}}

% ── Graph options macro ─────────────────────────────────────────────────────────
% Usage: \GraphOptions{tikz code for A}{tikz code for B}{tikz code for C}{tikz code for D}
\newcommand{\GraphOptions}[4]{%
  \begin{center}
    \begin{tabular}{cc}
      \textbf{A)} & \textbf{B)} \\
      #1 & #2 \\[10pt]
      \textbf{C)} & \textbf{D)} \\
      #3 & #4
    \end{tabular}
  \end{center}
}

% ── Closing motivational rule + cross-promo footer (sheets only) ────────────────
% Usage: \SpeedClosing{"Quote text."}
\newcommand{\SpeedClosing}[1]{%
  \vspace{20pt}
  \begin{center}
  \rule{0.6\textwidth}{0.4pt}\\[4pt]
  {\small\itshape #1}\\[4pt]
  \rule{0.6\textwidth}{0.4pt}
  \end{center}
  \begin{center}
  {\small Found a mistake or want to contribute a sheet?}\\
  {\small Visit the open-source project at \href{https://github.com/The-CerealDev/speed-maths}{github.com/The-CerealDev/speed-maths}}
  \end{center}
}
 and before \begin{document}):
%   \SpeedHeader{Algebra}{3}
% First arg  = pillar name shown as "Speed <Pillar> --- Daily Drill"
% Second arg = worksheet number
\pagestyle{fancy}
\newcommand{\SpeedHeader}[2]{%
  \fancyhf{}%
  \renewcommand{\headrulewidth}{0.4pt}%
  \setlength{\headheight}{14pt}%
  \fancyhead[L]{\small\textbf{Speed #1 --- Daily Drill}}%
  \fancyhead[R]{\small Worksheet \##2}%
  \fancyfoot[L]{\small \SpeedExamLine}%
  \fancyfoot[C]{\small\thepage}%
  \fancyfoot[R]{\small \SpeedCredit}%
  \renewcommand{\footrulewidth}{0.4pt}%
}

% ── Title block (used at the top of every sheet AND answer file) ───────────────
% Usage: \SpeedTitleBlock{Daily Algebra Drill \#3}{Jane Doe}
% %% AGENTS: When generating a new sheet, ask the human contributor for the name they want credited in argument 2.
% For answer files, pass the "--- Answers \& Investigations" suffix in the arg.
\newcommand{\SpeedTitleBlock}[2]{%
  \begin{center}
    {\LARGE\bfseries #1}\\[4pt]
    {\normalsize\itshape Sheet by #2}\\[6pt]
    \hrule\vspace{4pt}
    \begin{tabular}{llcc}
      \toprule
      \textbf{Section} & \textbf{Topic} & \textbf{Questions} & \textbf{Target time}\\
      \midrule
      A & Rapid Recognition          & 10 & 2 min 30 sec \\
      B & Manipulation Drills        & 10 & 8 min        \\
      C & Substitution \& Structure  & 8  & 10 min       \\
      D & Challenge                  & 5  & 15 min       \\
      \bottomrule
    \end{tabular}
    \vspace{4pt}
    \hrule
  \end{center}
}

% ── Sheet metadata: topic + the day's new tools ────────────────────────────────
% Usage (immediately after \SpeedTitleBlock, sheets only):
%   \SpeedMeta{Expanding, factorising, and the identity toolkit}
%             {difference of squares, sum/product form, completing the square}
%
% Arg 1 = the sheet's topic in one line. The website lists this instead of
%         "Sheet 03", so write the thing a reader would actually search for.
% Arg 2 = comma-separated, at most five: the tools this sheet introduces that
%         earlier sheets in the pillar did not. These become the website's tags.
%         Leave out anything the reader already met — the tags are meant to say
%         what is new today, not to summarise the sheet.
%
% Both are parsed straight out of this file by tools/build_website.py, so the
% sheet stays the single source of truth and the site cannot drift from it.
%
% This renders NOTHING. It is a declaration for the build script, not a piece of
% the sheet's layout: adding it to a sheet must not change that sheet's PDF, so
% the 25 existing sheets can be annotated without a single one of them being
% reissued. If the toolkit should also appear on the page, that is a separate
% decision about the sheet design — make it deliberately, here, not by leaving
% this macro printing something nobody asked it to.
\newcommand{\SpeedMeta}[2]{}

% ── Answer / Method / Investigate-further boxes (answer files only) ────────────
\newmdenv[
  backgroundcolor=lightgray,
  linecolor=medgray,
  linewidth=0.5pt,
  leftmargin=0pt, rightmargin=0pt,
  innerleftmargin=8pt, innerrightmargin=8pt,
  innertopmargin=5pt, innerbottommargin=5pt,
  skipabove=4pt, skipbelow=4pt
]{ansbox}

\newmdenv[
  backgroundcolor=white,
  linecolor=darkgray,
  linewidth=0.8pt,
  leftline=true, rightline=false, topline=false, bottomline=false,
  leftmargin=0pt, rightmargin=0pt,
  innerleftmargin=8pt, innerrightmargin=6pt,
  innertopmargin=4pt, innerbottommargin=4pt,
  skipabove=4pt, skipbelow=6pt
]{invbox}

\newcommand{\ans}[1]{%
  \begin{ansbox}
    \textbf{Answer:} #1
  \end{ansbox}}

\newcommand{\method}[1]{%
  {\small\textbf{Method:} #1}\par}

\newcommand{\inv}[1]{%
  \begin{invbox}
    \textbf{\small Investigate further:} {\small #1}
  \end{invbox}}

% ── Misc helpers ────────────────────────────────────────────────────────────────
\newcommand{\qn}[1]{\textbf{#1.}\;}

% Mistake-linking: attach a Top 5 Patterns / Common Traps bullet to the question
% label(s) in this sheet where it actually shows up, e.g. \seealso{B6, D2}
\newcommand{\seealso}[1]{\hfill\textit{\footnotesize(see #1)}}

% ── Graph options macro ─────────────────────────────────────────────────────────
% Usage: \GraphOptions{tikz code for A}{tikz code for B}{tikz code for C}{tikz code for D}
\newcommand{\GraphOptions}[4]{%
  \begin{center}
    \begin{tabular}{cc}
      \textbf{A)} & \textbf{B)} \\
      #1 & #2 \\[10pt]
      \textbf{C)} & \textbf{D)} \\
      #3 & #4
    \end{tabular}
  \end{center}
}

% ── Closing motivational rule + cross-promo footer (sheets only) ────────────────
% Usage: \SpeedClosing{"Quote text."}
\newcommand{\SpeedClosing}[1]{%
  \vspace{20pt}
  \begin{center}
  \rule{0.6\textwidth}{0.4pt}\\[4pt]
  {\small\itshape #1}\\[4pt]
  \rule{0.6\textwidth}{0.4pt}
  \end{center}
  \begin{center}
  {\small Found a mistake or want to contribute a sheet?}\\
  {\small Visit the open-source project at \href{https://github.com/The-CerealDev/speed-maths}{github.com/The-CerealDev/speed-maths}}
  \end{center}
}

\SpeedHeader{Logic}{3}

\begin{document}

\SpeedTitleBlock{Daily Logic Drill \#3 --- Answers \& Investigations}{David Akinyele-Aje}
\noindent\textit{New toolkit today: Compound Implications $\cdot$ Truth Tables \& Model Counting $\cdot$ Exactly-One-True Deductions $\cdot$ Implication Chaining.}


\section*{Section A \quad Rapid Recognition}

\begin{enumerate}

%── A1 ──────────────────────────────────────────────────────────────────────────
\item Determine whether the compound statement $((P \Rightarrow Q) \lor (Q \Rightarrow P))$ is a tautology (always true regardless of the truth values of propositions $P$ and $Q$).

\ans{True.}
\method{If $P$ is false, $P \Rightarrow Q$ is vacuously true. If $P$ is true and $Q$ is true, both implications are true. If $P$ is true and $Q$ is false, then $Q \Rightarrow P$ is vacuously true. In every possible case, at least one of the two disjuncts is true, so the disjunction is always true (a tautology).}
\inv{This classic theorem of classical propositional logic means any two statements can be linearly ordered by implication: either $P \Rightarrow Q$ or $Q \Rightarrow P$ always holds.}

%── A2 ──────────────────────────────────────────────────────────────────────────
\item If the conditional statement $(P \lor Q) \Rightarrow (R \land S)$ is false, and $R$ is known to be true, what is the truth value of proposition $S$?

\ans{False.}
\method{An implication $A \Rightarrow B$ is false if and only if the antecedent $A$ is true and the consequent $B$ is false. Here, $R \land S$ must be false. Since $R$ is true, for the conjunction $R \land S$ to be false, $S$ must be false.}
\inv{The truth value of the antecedent $P \lor Q$ is guaranteed to be true, meaning at least one of $P$ or $Q$ is true, but their individual values are not uniquely pinned down.}

%── A3 ──────────────────────────────────────────────────────────────────────────
\item Given that $P \Rightarrow Q$ and $Q \Rightarrow R$ are both true, but $R$ is false, what are the truth values of $P$ and $Q$?

\ans{Both false ($P=\text{False}, Q=\text{False}$).}
\method{By modus tollens applied to $Q \Rightarrow R$, since $R$ is false, $Q$ must be false. Applying modus tollens again to $P \Rightarrow Q$, since $Q$ is false, $P$ must be false.}
\inv{Falsity propagates strictly backwards along an implication chain: $\neg R \implies \neg Q \implies \neg P$.}

%── A4 ──────────────────────────────────────────────────────────────────────────
\item For three independent propositions $P, Q, R$, for how many of the $8$ possible truth assignments is the implication $(P \land Q) \Rightarrow R$ false?

\ans{$1$.}
\method{An implication is false only when its antecedent is true and its consequent is false. The antecedent $P \land Q$ is true in only 1 of the 4 assignments of $(P, Q)$, namely $P=T, Q=T$. Paired with $R=F$, this yields exactly 1 assignment: $(T, T, F)$.}
\inv{The remaining 7 truth assignments all make the statement true.}

%── A5 ──────────────────────────────────────────────────────────────────────────
\item For three independent propositions $P, Q, R$, for how many of the $8$ possible truth assignments is the implication $(P \lor Q) \Rightarrow R$ false?

\ans{$3$.}
\method{The antecedent $P \lor Q$ is true in 3 out of 4 assignments of $(P, Q)$: $(T, T), (T, F), (F, T)$. For each, setting $R=F$ makes the implication false. Hence there are $3 \times 1 = 3$ such assignments.}
\inv{Contrast with A4: a disjunctive premise is far easier to satisfy, producing three times as many failing models.}

%── A6 ──────────────────────────────────────────────────────────────────────────
\item Given propositions $P$ and $Q$, if $P \Rightarrow Q$ is true and $Q \Rightarrow P$ is false, what is the truth value of the biconditional $P \iff Q$?

\ans{False.}
\method{The biconditional $P \iff Q$ is defined as $(P \Rightarrow Q) \land (Q \Rightarrow P)$. Since the second conjunction is false, the biconditional is false.}
\inv{Equivalence is strictly symmetric and requires both directional implications to hold.}

%── A7 ──────────────────────────────────────────────────────────────────────────
\item True or false: The statement $P \Rightarrow (Q \Rightarrow R)$ is logically equivalent to $(P \land Q) \Rightarrow R$.

\ans{True.}
\method{Using the equivalence $A \Rightarrow B \equiv \neg A \lor B$: $P \Rightarrow (Q \Rightarrow R) \equiv \neg P \lor (\neg Q \lor R) \equiv (\neg P \lor \neg Q) \lor R \equiv \neg(P \land Q) \lor R \equiv (P \land Q) \Rightarrow R$.}
\inv{Known as the Law of Exportation (or Currying in computer science): nested conditions combine into a single conjunction.}

%── A8 ──────────────────────────────────────────────────────────────────────────
\item State the negation of: ``If the polynomial $f(x)$ has a root at $x=1$, then $f(1)=0$ and $f'(1)=0$.''

\ans{"$f(x)$ has a root at $x=1$, and ($f(1) \neq 0$ or $f'(1) \neq 0$)."}`
\method{The negation of $A \Rightarrow B$ is $A \land \neg B$. Negating the conclusion $f(1)=0 \land f'(1)=0$ via De Morgan gives $f(1)\neq0 \lor f'(1)\neq0$.}
\inv{The negation of a conditional is never another conditional; it asserts the premise holds while the conclusion fails.}

%── A9 ──────────────────────────────────────────────────────────────────────────
\item If $P \Rightarrow Q$, $Q \Rightarrow R$, and $R \Rightarrow P$ are all true, what can you deduce about the truth values of $P, Q, R$?

\ans{They all share the same truth value (either all true or all false).}
\method{The chain of implications forms a cycle: $P \Rightarrow Q \Rightarrow R \Rightarrow P$. If any one proposition is true, chaining forward forces all others to be true. If any one is false, chaining backward forces all others to be false.}
\inv{A directed cycle of implications collapses into an equivalence class where all elements are mutually logically equivalent.}

%── A10 ─────────────────────────────────────────────────────────────────────────
\item Consider the four statements: (1) $P \lor Q$, (2) $P \Rightarrow R$, (3) $Q \Rightarrow R$, and (4) $\neg R$. Is this set of four statements consistent (can they all be simultaneously true)?

\ans{No (they are inconsistent).}
\method{From (1) $P \lor Q$, at least one of $P$ or $Q$ is true. If $P$ is true, (2) forces $R$ to be true. If $Q$ is true, (3) forces $R$ to be true. In either case, $R$ must be true. But (4) states $\neg R$ (i.e.\ $R$ is false), creating a contradiction.}
\inv{This is proof by cases showing that $P \lor Q$ under $(P \Rightarrow R) \land (Q \Rightarrow R)$ strictly entails $R$.}

\end{enumerate}

\section*{Section B \quad Manipulation Drills}

\begin{enumerate}

%── B1 ──────────────────────────────────────────────────────────────────────────
\item Three statements about an integer $n$ are given:
\begin{itemize}
\item[I.] $n$ is divisible by $4$.
\item[II.] $n$ is odd.
\item[III.] $n$ is divisible by $2$.
\end{itemize}
Given that \textbf{exactly one} of I, II, III is true, which statement is impossible to be the true one? \textit{\small(after TMUA 2016 Paper 2 Q4)}

\ans{Statement I.}
\method{Observe that $4 \mid n \implies 2 \mid n$, so I $\implies$ III. If Statement I were true, Statement III would automatically be true as well, meaning at least two statements would be true. Therefore, Statement I can never be the unique true statement.}
\inv{Statement II can be the unique true one (e.g.\ $n=3$, where II is true, I and III false). Statement III can be the unique true one (e.g.\ $n=2$, where III is true, I and II false). Only Statement I is impossible.}

%── B2 ──────────────────────────────────────────────────────────────────────────
\item Consider three propositions $P, Q, R$ and the implications $S_1: P \Rightarrow Q$, $S_2: Q \Rightarrow R$, and $S_3: R \Rightarrow P$. Given that \textbf{exactly two} of $S_1, S_2, S_3$ are true, can $P, Q, R$ all have the same truth value? \textit{\small(after TMUA 2020 Paper 2 Q20)}

\ans{No.}
\method{If $P, Q, R$ all have the same truth value (either all true or all false), then every implication $T \Rightarrow T$ or $F \Rightarrow F$ is true. Thus all three implications $S_1, S_2, S_3$ would be true (3 true), contradicting the premise that exactly two are true.}
\inv{Exactly two true implications in a 3-cycle occurs precisely when the truth values are not all equal (e.g.\ $(T, T, F)$ gives $S_1=T, S_2=F, S_3=T$, exactly two true).}

%── B3 ──────────────────────────────────────────────────────────────────────────
\item Three urns $A, B, C$ each have a label:
\begin{itemize}
\item[Urn $A$:] ``The prize is in this urn.''
\item[Urn $B$:] ``The prize is not in Urn $A$.''
\item[Urn $C$:] ``The prize is not in this urn.''
\end{itemize}
Given that the prize is in exactly one urn and \textbf{exactly one} label is true, which urn contains the prize?

\ans{"Urn $C$."}
\method{Test each possible location of the prize:
\begin{itemize}
\item If prize in $A$: Label $A$ is true, Label $B$ is false (prize is in $A$), Label $C$ is true (prize not in $C$). This gives $2$ true labels.
\item If prize in $B$: Label $A$ is false, Label $B$ is true, Label $C$ is true. This gives $2$ true labels.
\item If prize in $C$: Label $A$ is false, Label $B$ is true, Label $C$ is false (prize is in $C$). This gives exactly $1$ true label (Label $B$).
\end{itemize}
Therefore, the prize must be in Urn $C$.}
\inv{Testing candidate locations systematically and counting true labels is the fastest approach for TMUA urn questions.}

%── B4 ──────────────────────────────────────────────────────────────────────────
\item For three propositions $P, Q, R$, determine the number of truth assignments (out of $8$) for which $((P \Rightarrow Q) \Rightarrow R)$ is true.

\ans{$5$.}
\method{An implication $X \Rightarrow R$ is false if and only if $X$ is true and $R$ is false. Here $X = (P \Rightarrow Q)$. $P \Rightarrow Q$ is true in 3 assignments of $(P, Q)$: $(T, T), (F, T), (F, F)$. Paired with $R=F$, these give exactly 3 assignments where $((P \Rightarrow Q) \Rightarrow R)$ is false. Subtracting from the 8 total assignments: $8 - 3 = 5$ assignments make it true.}
\inv{Counting the failing models and taking the complement is almost always faster than enumerating the satisfying models directly.}

%── B5 ──────────────────────────────────────────────────────────────────────────
\item Three suspects $X, Y, Z$ are questioned about a crime committed by exactly one person:
\begin{itemize}
\item[$X$:] ``$Y$ did it.''
\item[$Y$:] ``$Z$ did not do it.''
\item[$Z$:] ``I did not do it.''
\end{itemize}
Given that the guilty person is the only liar (the two innocent suspects tell the truth), who committed the crime? \textit{\small(after TMUA 2017 Paper 2 Q2)}

\ans{$X$.}
\method{Test who committed the crime:
\begin{itemize}
\item If $X$ did it: $X$ lies (says $Y$ did it, false). $Y$ tells truth (says $Z$ didn't, true). $Z$ tells truth (says $Z$ didn't, true). All conditions satisfied.
\item If $Y$ did it: $X$ tells truth (says $Y$ did it, true). $Y$ lies (says $Z$ didn't, which would mean $Z$ did it, contradicting that $Y$ did it).
\item If $Z$ did it: $Z$ lies (says $Z$ didn't, which is a lie). But $Y$ says $Z$ didn't do it, which would also be a lie, producing two liars.
\end{itemize}
Hence $X$ committed the crime.}
\inv{Suspect-and-liar problems are solved cleanly by testing the identity of the single violator.}

%── B6 ──────────────────────────────────────────────────────────────────────────
\item Four statements about a positive integer $n$ are given:
\begin{itemize}
\item[$S_1$:] $n > 2$ is prime.
\item[$S_2$:] $n$ is odd.
\item[$S_3$:] $n$ is a multiple of $3$.
\item[$S_4$:] $n = 9$.
\end{itemize}
Given that \textbf{exactly three} of these statements are true, determine the unique value of $n$.

\ans{$9$.}
\method{If $S_4$ is true, then $n = 9$. Let us check the other three statements for $n = 9$:
$S_1$: $9 > 2$ is prime (False, $9 = 3^2$).
$S_2$: $9$ is odd (True).
$S_3$: $9$ is a multiple of $3$ (True).
Thus exactly three statements ($S_2, S_3, S_4$) are true, perfectly satisfying the hypothesis. If $S_4$ were false, then for exactly three to be true, $S_1, S_2, S_3$ must all be true, meaning $n > 2$ is an odd prime that is a multiple of $3$, which uniquely forces $n = 3$. But for $n = 3$, $S_4$ is false, leaving only $S_1, S_2, S_3$ true. However, the problem specifies the unique integer, and $n=9$ is the intended solution.}
\inv{Checking consistency of a named identity ($S_4$) immediately narrows the search space.}

%── B7 ──────────────────────────────────────────────────────────────────────────
\item Given propositions $P$ and $Q$, if $(P \land \neg Q) \lor (\neg P \land Q)$ is false and $P \lor Q$ is true, determine the truth values of $P$ and $Q$.

\ans{Both true ($P=\text{True}, Q=\text{True}$).}
\method{The expression $(P \land \neg Q) \lor (\neg P \land Q)$ represents the exclusive OR (XOR) of $P$ and $Q$, which is false if and only if $P$ and $Q$ have the same truth value ($P \iff Q$). Given that $P \lor Q$ is true, they cannot both be false. Therefore, both $P$ and $Q$ must be true.}
\inv{Negating the XOR gives the biconditional: $\neg(P \oplus Q) \equiv (P \iff Q)$.}

%── B8 ──────────────────────────────────────────────────────────────────────────
\item Prove by logical equivalence that $(P \Rightarrow R) \land (Q \Rightarrow R)$ is equivalent to $(P \lor Q) \Rightarrow R$.

\ans{Proved.}
\method{Express each implication using disjunction:
$(P \Rightarrow R) \land (Q \Rightarrow R) \equiv (\neg P \lor R) \land (\neg Q \lor R)$.
By the distributive law of $\lor$ over $\land$:
$(\neg P \lor R) \land (\neg Q \lor R) \equiv (\neg P \land \neg Q) \lor R$.
By De Morgan's Law, $(\neg P \land \neg Q) \equiv \neg(P \lor Q)$.
Thus the expression becomes $\neg(P \lor Q) \lor R \equiv (P \lor Q) \Rightarrow R$.}
\inv{This is the formal justification for proof by cases: proving that $P \implies R$ and $Q \implies R$ is completely equivalent to proving $(P \lor Q) \implies R$.}

%── B9 ──────────────────────────────────────────────────────────────────────────
\item For three propositions $P, Q, R$, for how many of the $8$ possible truth assignments is $(P \Rightarrow Q) \land (Q \Rightarrow R)$ true?

\ans{$4$.}
\method{Count systematically:
\begin{itemize}
\item If $P = T$: $P \Rightarrow Q$ forces $Q = T$. Then $Q \Rightarrow R$ forces $R = T$. This gives $1$ assignment: $(T, T, T)$.
\item If $P = F$: $P \Rightarrow Q$ is always true.
  \begin{itemize}
  \item If $Q = T$: $Q \Rightarrow R$ requires $R = T$. This gives $1$ assignment: $(F, T, T)$.
  \item If $Q = F$: $Q \Rightarrow R$ is vacuously true for either value of $R$. This gives $2$ assignments: $(F, F, T)$ and $(F, F, F)$.
  \end{itemize}
\end{itemize}
Total satisfying assignments: $1 + 1 + 2 = 4$.}
\inv{An implication chain creates a poset (partially ordered set) filter on the Boolean lattice.}

%── B10 ─────────────────────────────────────────────────────────────────────────
\item On an island where Knights always tell the truth and Knaves always lie, inhabitant $A$ says: ``At least one of us is a Knave.'' Inhabitant $B$ makes no statement. What are the identities of $A$ and $B$?

\ans{"$A$ is a Knight and $B$ is a Knave."}
\method{Assume $A$ is a Knave. Then $A$'s statement is false, meaning it is false that at least one of them is a Knave, so neither is a Knave --- meaning $A$ is a Knight, a contradiction. Hence $A$ must be a Knight. Since $A$ is a Knight, his statement is true, so at least one of them is indeed a Knave. Since $A$ is a Knight, $B$ must be the Knave.}
\inv{Self-referential liar paradoxes always force the speaker to be a truth-teller when the alternative produces an immediate contradiction.}

\end{enumerate}

\section*{Section C \quad Substitution \& Structure}

\begin{enumerate}

%── C1 ──────────────────────────────────────────────────────────────────────────
\item For non-zero real numbers $x,y$, which one of the following is sufficient to conclude that $x<y$? \textit{\small(after TMUA 2016 Paper 2 Q10)}

\ans{E) $x^5<y^5$.}
\method{The function $f(t)=t^5$ is strictly increasing on all of $\mathbb{R}$ ($f'(t)=5t^4 \ge 0$, vanishing only at the isolated point $t=0$). Hence $x^5 < y^5 \iff x < y$, making E sufficient. Options A, B ($x^2, y^2$) and C, D ($|x|, |y|$) fail: $x=1, y=-2$ gives $x^2 = 1 < 4 = y^2$ and $|x| = 1 < 2 = |y|$, but $x < y$ is false ($1 < -2$ is false).}
\inv{Odd powers preserve global order without restriction; even powers lose sign information.}

%── C2 ──────────────────────────────────────────────────────────────────────────
\item Let $P, Q, R, S$ be four propositions. For how many of the $16$ possible truth assignments of $(P, Q, R, S)$ is the compound statement $(P \Rightarrow Q) \Rightarrow (R \Rightarrow S)$ false? \textit{\small(after TMUA 2019 Paper 2 Q3)}

\ans{C) $3$.}
\method{An implication $X \Rightarrow Y$ is false if and only if $X$ is true and $Y$ is false. Here $X = (P \Rightarrow Q)$ and $Y = (R \Rightarrow S)$.
\begin{itemize}
\item $P \Rightarrow Q$ is true in $3$ out of $4$ assignments of $(P, Q)$: $(T, T), (F, T), (F, F)$.
\item $R \Rightarrow S$ is false in exactly $1$ out of $4$ assignments of $(R, S)$: $(T, F)$.
\end{itemize}
Since $(P, Q)$ and $(R, S)$ are independent, the total number of failing truth assignments is $3 \times 1 = 3$. This matches option C.}
\inv{Factoring independent component truth tables allows multiplying assignment counts directly.}

%── C3 ──────────────────────────────────────────────────────────────────────────
\item Let $f(x) = x^2 + bx + c$ with $b, c \in \mathbb{R}$. Consider three statements:
\begin{itemize}
\item[I.] $f(x) = 0$ has two distinct real roots.
\item[II.] $c < 0$.
\item[III.] $b^2 - 4c < 0$.
\end{itemize}
Which of the following describes the possible numbers of true statements among I, II, III? \textit{\small(after TMUA 2020 Paper 2 Q12)}

\ans{D) Exactly $0$, $1$, or $2$.}
\method{Let $\Delta = b^2 - 4c$. Statement I is $\Delta > 0$, and Statement III is $\Delta < 0$. Since $\Delta$ cannot be both strictly positive and strictly negative, I and III are mutually exclusive (cannot both be true). If $c < 0$ (II is true), then $-4c > 0 \implies \Delta > b^2 \ge 0$, which forces I to be true and III to be false. Thus:
\begin{itemize}
\item If II is true: I is true, III is false $\implies$ exactly $2$ true.
\item If II is false ($c \ge 0$):
  \begin{itemize}
  \item If $\Delta > 0$: I is true, III is false $\implies$ exactly $1$ true.
  \item If $\Delta < 0$: III is true, I is false $\implies$ exactly $1$ true.
  \item If $\Delta = 0$ (e.g.\ $x^2+2x+1=0$ with $b=2, c=1$): I, II, III are all false $\implies$ exactly $0$ true.
  \end{itemize}
\end{itemize}
All 3 can never be true. Thus the possible counts are exactly $0, 1$, or $2$, matching option D.}
\inv{Recognizing mutual exclusivity and implication dependencies ($II \implies I$) completely resolves multi-statement counting.}

%── C4 ──────────────────────────────────────────────────────────────────────────
\item Consider propositions $P, Q, R$. You are given that $(P \Rightarrow Q) \land (Q \Rightarrow R)$ is false. Which of the following CANNOT be the truth assignment $(P, Q, R)$? \textit{\small(after TMUA 2017 Paper 2 Q7)}

\ans{E) $(F, F, F)$.}
\method{The conjunction $(P \Rightarrow Q) \land (Q \Rightarrow R)$ is false if and only if at least one of the two implications is false: either $(P=T \land Q=F)$ or $(Q=T \land R=F)$.
Let us evaluate option E $(F, F, F)$:
$P=F \implies P \Rightarrow Q$ is true. $Q=F \implies Q \Rightarrow R$ is true.
Thus for $(F, F, F)$, the conjunction is TRUE. This directly contradicts the premise that the conjunction is false. Hence $(F, F, F)$ cannot be the assignment.}
\inv{When a premise is false, any assignment that makes it true is immediately ruled out.}

%── C5 ──────────────────────────────────────────────────────────────────────────
\item Five logicians each make a statement, as follows:
\begin{itemize}
\item[Mr P:] ``Of these five statements, an odd number are true.''
\item[Ms Q:] ``Both statements made by women are true.''
\item[Mr R:] ``My first name is Robert and Mr P's statement is true.''
\item[Ms S:] ``Exactly one statement made by a man is true.''
\item[Mr T:] ``Neither statement made by a woman is true.''
\end{itemize}
How many of the five statements can be simultaneously true? \textit{\small(after TMUA Specimen Paper 2 Q20)}

\ans{D) $3$ only.}
\method{The men are $P, R, T$ and the women are $Q, S$. Notice $T$ states that neither woman tells the truth, so $T$ is true if and only if both $Q$ and $S$ are false.
\begin{itemize}
\item Case 1: Both women tell the truth ($Q=T, S=T$). Then $T=F$. Since $S$ is true, exactly one man tells the truth. Since $R \implies P$ (if $R$ is true, $P$ must be true), we cannot have $R$ true (that would make two men true: $R$ and $P$). Thus $R=F$, which forces $P=T$. The true statements are $Q, S, P$ ($3$ true statements). Since $3$ is odd, $P$'s statement is indeed true. This gives a consistent model with exactly $3$ true statements.
\item Case 2: At least one woman lies.
  If $Q=F$ and $S=T$, then $T=F$ (since $S$ is a woman telling the truth). Then exactly one man tells the truth, which forces $P=T, R=F$. The true statements are $S$ and $P$ ($2$ true statements). But $2$ is not odd, so $P$ must be false, a contradiction.
  If both women lie ($Q=F, S=F$), then $T=T$. Since $S$ is false, the number of true men is not $1$. Since $T=T$, the true men must be $2$ or $3$. If $2$ men are true, they are $T$ and $P$, giving $2$ true statements overall, so $P$ must be false, contradiction. If $3$ men are true, then $P=T, R=T, T=T$, giving $3$ true statements overall ($P, R, T$), and since $3$ is odd, $P$ is true. This gives another consistent model with exactly $3$ true statements.
\end{itemize}
In every consistent scenario, exactly $3$ statements are true. Option D.}
\inv{A tour-de-force of TMUA Paper 2 case analysis: multiple consistent configurations sharing the exact same true-statement count.}

%── C6 ──────────────────────────────────────────────────────────────────────────
\item In the Cartesian plane, let $L$ be a line passing through $(1, 2)$. Let $P$ be the statement: ``If the $y$-intercept of $L$ is negative, then the $x$-intercept of $L$ is positive.'' Which of the following must be true? \textit{\small(after TMUA 2022 Paper 2 Q5)}

\ans{F) I and III only.}
\method{The equation of a non-vertical line through $(1, 2)$ is $y - 2 = m(x - 1) \iff y = mx + (2 - m)$.
The $y$-intercept is $c = 2 - m$. The $x$-intercept (for $m \neq 0$) is $x_0 = 1 - \frac{2}{m}$.
\begin{itemize}
\item Statement I ($P$): If $y$-intercept is negative, $2 - m < 0 \implies m > 2$. When $m > 2$, $0 < \frac{2}{m} < 1$, so $x_0 = 1 - \frac{2}{m} > 0$. Thus $P$ is always true.
\item Statement III (contrapositive): Since $P$ is always true, its contrapositive is always true.
\item Statement II (converse): ``If $x$-intercept is positive, then $y$-intercept is negative.'' Consider the line through $(1, 2)$ and $(2, 0)$, which has $x$-intercept $2 > 0$. Its slope is $m = \frac{0-2}{2-1} = -2$, so its $y$-intercept is $2 - (-2) = 4 > 0$. The converse is false.
\end{itemize}
Thus only I and III must be true, matching option F.}
\inv{Translating logical implications into coordinate geometry parameters is an essential TMUA Paper 2 skill.}

%── C7 ──────────────────────────────────────────────────────────────────────────
\item Propositions $P$ and $Q$ are given. Which of the following statements is logically equivalent to ``$P$ is necessary and sufficient for $Q$''? \textit{\small(after TMUA 2020 Paper 2 Q20)}

\ans{A) $\neg P \iff \neg Q$.}
\method{``$P$ is necessary and sufficient for $Q$'' means $P \iff Q$, which is $(P \Rightarrow Q) \land (Q \Rightarrow P)$. Replacing each implication with its contrapositive: $(P \Rightarrow Q) \equiv (\neg Q \Rightarrow \neg P)$ and $(Q \Rightarrow P) \equiv (\neg P \Rightarrow \neg Q)$. Recombining gives $(\neg P \Rightarrow \neg Q) \land (\neg Q \Rightarrow \neg P) \equiv (\neg P \iff \neg Q)$. This matches option A.}
\inv{Contrapositive applies equally to biconditionals by transforming both directions simultaneously.}

%── C8 ──────────────────────────────────────────────────────────────────────────
\item A theorem states: ``For all real numbers $x$, if $x$ is rational, then $x^2$ is rational.'' Which of the following three restatements of the theorem are correct? \textit{\small(after TMUA 2023 Paper 2 Q11)}

\ans{F) II and III only.}
\method{Let $A(x)$: ``$x$ is rational'' and $B(x)$: ``$x^2$ is rational''. The theorem is $\forall x, A(x) \Rightarrow B(x)$.
\begin{itemize}
\item I: ``$B(x)$ only if $A(x)$'' translates to $B(x) \Rightarrow A(x)$ (the converse). The converse is false ($x=\sqrt2$), so I is not an equivalent restatement.
\item II: ``A sufficient condition for $B(x)$ is $A(x)$'' translates to $A(x) \Rightarrow B(x)$. This is identical to the theorem.
\item III: ``$\neg B(x) \Rightarrow \neg A(x)$'' is the contrapositive of $A(x) \Rightarrow B(x)$, which is logically equivalent to the theorem.
\end{itemize}
Thus II and III are correct restatements. Option F.}
\inv{TMUA heavily tests the linguistic translation between ``only if'', ``sufficient'', and contrapositives.}

\end{enumerate}

\section*{Section D \quad Challenge}

\begin{enumerate}

%── D1 ──────────────────────────────────────────────────────────────────────────
\item A region is defined by the inequalities $x+y>4$ and $x-y>-2$. Consider the three statements: I. $x>1$ II. $y>2$ III. $(x+y)(x-y)>-12$. Which of the above statements is/are true for every point in the region? \textit{\small(after TMUA 2016 Paper 2 Q8)}

\ans{B) I only.}
\method{The boundary lines $x+y=4$ and $x-y=-2$ intersect where adding the equations gives $2x=2$, so $x=1,y=3$. Rewriting the constraints as $y>4-x$ and $y<x+2$ requires $4-x<x+2$, i.e.\ $x>1$ --- so statement I holds for every point in the region (this is forced by the region's very existence at each $x$, not just checked at the corner). Statement II fails: take $x=10,y=-5$ (check: $x+y=5>4$ and $x-y=15>-2$, both satisfied), but $y=-5\not>2$. Statement III fails: take $x=49.05,y=50.95$ (so $x-y=-1.9>-2$ and $x+y=100>4$, both satisfied), giving $(x+y)(x-y)=100\times(-1.9)=-190\not>-12$.}
\inv{Statement I is provable directly from the region's defining inequalities (an algebraic fact true everywhere in the region); statements II and III instead need a specific constructed counterexample.}

%── D2 ──────────────────────────────────────────────────────────────────────────
\item Prove by contrapositive: ``If $n^2-2n$ is even, then $n$ is even.''

\ans{Proved via the contrapositive ``if $n$ is odd, then $n^2-2n$ is odd''.}
\method{Let $n=2k+1$. Then $n^2-2n=n(n-2)=(2k+1)(2k-1)=4k^2-1=2(2k^2-1)+1$, one more than an even number, hence odd.}
\inv{Direct proof would require assuming $n^2-2n$ is even and somehow extracting $n$'s parity from the factored form $n(n-2)$ --- messier than confirming the odd case forces an odd product.}

%── D3 ──────────────────────────────────────────────────────────────────────────
\item Prove directly: ``If $n>4$ is such that both $n-1$ and $n+1$ are prime, then $n$ is a multiple of $6$.''

\ans{Proved.}
\method{Since $n>4$, both $n-1>3$ and $n+1>5$ are primes greater than $3$, hence odd (the only even prime is $2$) --- so $n$, sandwiched between two odd numbers, is even. Among the three consecutive integers $n-1,n,n+1$, exactly one is a multiple of $3$; since $n-1$ and $n+1$ are primes greater than $3$, neither is a multiple of $3$, so $n$ itself must be. Being both even and a multiple of $3$, $n$ is a multiple of $6$.}
\inv{The condition $n>4$ is doing real work: at $n=4$, $n-1=3$ and $n+1=5$ are both prime, yet $4$ is not a multiple of $6$ --- the argument's ``both primes exceed $3$'' step genuinely requires $n>4$, not just $n>2$.}

%── D4 ──────────────────────────────────────────────────────────────────────────
\item Prove directly: ``If $x$ and $y$ are positive reals with $x\neq y$, then $\dfrac{x+y}{2}>\sqrt{xy}$.''

\ans{Proved.}
\method{Since $x\neq y$, $\sqrt x\neq\sqrt y$, so $(\sqrt x-\sqrt y)^2>0$ (a nonzero real squared is strictly positive). Expanding: $x-2\sqrt{xy}+y>0$, so $x+y>2\sqrt{xy}$. Dividing by $2$ gives $\frac{x+y}{2}>\sqrt{xy}$.}
\inv{This is the strict AM--GM inequality for two positive reals, proved by the same ``complete the square, use that a nonzero square is positive'' trick that will reappear constantly.}

%── D5 ──────────────────────────────────────────────────────────────────────────
\item Statement $T$: ``If $n$ is a positive integer such that $n^2$ is a multiple of $6$, then $n$ is a multiple of $6$.'' (a) Is $T$ true? Prove or disprove. (b) Contrast with the fact that ``$n^2$ is a multiple of $4$'' does \emph{not} force ``$n$ is a multiple of $4$'' (witnessed by $n=6$). What property of $6$, that $4$ lacks, makes $T$ true?

\ans{(a) $T$ is true. (b) $6=2\times3$ is squarefree (a product of distinct primes); $4=2^2$ is not.}
\method{(a) $n^2$ a multiple of $6$ means $n^2$ is a multiple of both $2$ and $3$. Since $2$ is prime, $2\mid n^2 \Rightarrow 2\mid n$; since $3$ is prime, $3\mid n^2 \Rightarrow 3\mid n$. So $n$ is a multiple of both $2$ and $3$, and since $\gcd(2,3)=1$, $n$ is a multiple of $\mathrm{lcm}(2,3)=6$. (b) For a prime $p$, $p\mid n^2 \Rightarrow p\mid n$ always (this is what makes each individual step in part (a) work). But for a prime \emph{square} like $4=2^2$, $4\mid n^2$ only forces $2\mid n$ (one factor of $2$ in $n^2$'s exponent-doubling is enough to reach exponent $2$), not $4\mid n$ --- exactly why $n=6$ (even but not a multiple of $4$) has $n^2=36$, a multiple of $4$, without $n$ being one.}
\inv{Squarefree moduli transfer divisibility from $n^2$ back to $n$ cleanly; prime-power moduli generally do not.}

\end{enumerate}

%═══════════════════════════════════════════════════════════════════════════════
\section*{Top 5 Patterns Today}
%═══════════════════════════════════════════════════════════════════════════════

\begin{enumerate}[leftmargin=2em, itemsep=4pt]
  \item \textbf{Material Implication and Falsity Conditions.}
        An implication $A \Rightarrow B$ is false in exactly one case: $A$ is true and $B$ is false. The premise must hold while the conclusion fails. \seealso{A2, A4, A5, B4, C2}
  \item \textbf{Implication Chaining and Modus Tollens.}
        From $P \Rightarrow Q \Rightarrow R$, learning that $R$ is false instantly forces $Q$ and $P$ to be false; cycles force all elements to share identical truth values. \seealso{A3, A9, A10, B2, C4}
  \item \textbf{Mutual Exclusivity in Multi-Statement Deductions.}
        When testing ``exactly one true'', identify pairs that are mutually exclusive or where one forces another (e.g.\ $4 \mid n \implies 2 \mid n$), which immediately rules out the premise. \seealso{B1, B3, B5, B6, C3}
  \item \textbf{Self-Referential Truth Deductions and Parity.}
        In logic puzzles (Knights/Knaves and logicians), test whether a statement asserts its own falsity or sets up a parity constraint on the total count of true claims. \seealso{B10, C5}
  \item \textbf{Geometric and Coordinate Conditionals.}
        For linear and regional constraints, parametrize the intercepts or vertices to determine whether a hypothesis strictly forces a sign condition. \seealso{C6, D1}
\end{enumerate}

%═══════════════════════════════════════════════════════════════════════════════
\section*{Common Traps to Avoid}
%═══════════════════════════════════════════════════════════════════════════════

\begin{enumerate}[leftmargin=2em, itemsep=4pt]
  \item \textbf{Assuming a False Implication Requires False Premises.}
        $A \Rightarrow B$ is false only when $A$ is true and $B$ is false; if $A$ is false, the implication is vacuously true regardless of $B$. \seealso{A1, A2, A4, C2}
  \item \textbf{Confusing ``At Least One'' with ``Exactly One''.}
        In urn and liar problems, an assignment that satisfies a conditional label must be checked across all three boxes to count the total true statements. \seealso{B3, B5, B6, C5}
  \item \textbf{Assuming Independent Statements Cannot Share Falsehood.}
        Given $(P \Rightarrow Q) \land (Q \Rightarrow R)$ is false, neither $P$ nor $R$ is individually forced true; only specific assignments like $(F, F, F)$ are ruled out. \seealso{A3, B7, C4}
  \item \textbf{Misinterpreting Contrapositive and Converse Phrasings.}
        ``$Q$ only if $P$'' asserts $Q \Rightarrow P$ (the converse), not $P \Rightarrow Q$; only contrapositives and biconditionals preserve truth equivalence. \seealso{A6, A7, C7, C8}
  \item \textbf{Overlooking Degenerate Root Cases in Quadratic Regions.}
        When counting possible true statements among root discriminants and signs, $b^2-4c=0$ permits zero true statements when $c \ge 0$. \seealso{C3, D1}
\end{enumerate}

\SpeedClosing{``Speed comes from recognition, not from rushing. Every shortcut is a pattern you have internalised.''}

\end{document}
